\documentclass{article}
\usepackage{amsmath,amssymb,amsthm}
\usepackage{algorithm,algorithmic}
\usepackage{hyperref}
\usepackage{enumitem}
\newtheorem{theorem}{Theorem}
\newtheorem{lemma}{Lemma}
\newtheorem{assumption}{Assumption}
\newcommand{\norm}[1]{\left\lVert#1\right\rVert}
\begin{document}

\section*{AdaGrad (Scalar Version)}

\section*{Mathematical Formulation}
Let $f:\mathbb{R}^{d}\to\mathbb{R}$ be a convex and $L$‑smooth objective.  
Denote by $g_{k}= \nabla f(w_{k})\in\mathbb{R}^{d}$ the (full) gradient at iterate $w_{k}\in\mathbb{R}^{d}$.  
Define the accumulated squared gradient
\[
G_{k}\;:=\;\sum_{i=1}^{k} \|g_{i}\|^{2},\qquad G_{0}=0,
\]
and the adaptive stepsize
\[
\eta_{k}\;:=\;\frac{\alpha}{\sqrt{G_{k}}},\qquad \alpha>0\text{ is a user‑chosen constant.}
\]
The update rule is
\[
\boxed{w_{k+1}=w_{k}-\eta_{k}\,g_{k}
      =w_{k}-\frac{\alpha}{\sqrt{G_{k}}}\,g_{k}}.
\]
For notational convenience we also keep the running sum of raw gradients
\[
A_{k}= \sum_{i=1}^{k} g_{i},\qquad A_{0}=0,
\]
although the algorithm never uses $A_{k}$ directly.

\section*{Pseudocode}
\begin{algorithm}[H]
\caption{Scalar AdaGrad}
\begin{algorithmic}[1]
\STATE \textbf{Input:} initial point $w_{0}\in\mathbb{R}^{d}$, 
        learning‑rate constant $\alpha>0$, number of iterations $T$.
\STATE $G_{0}\gets 0$ \COMMENT{accumulated squared gradient}
\FOR{$k=0$ {\bf to} $T-1$}
    \STATE $g_{k}\gets \nabla f(w_{k})$
    \STATE $G_{k+1}\gets G_{k}+ \|g_{k}\|^{2}$
    \STATE $\eta_{k}\gets \alpha/\sqrt{G_{k+1}}$
    \STATE $w_{k+1}\gets w_{k}-\eta_{k}\,g_{k}$
\ENDFOR
\STATE \textbf{Output:} $\bar w_{T}= \frac{1}{T}\sum_{k=1}^{T} w_{k}$ (optional average)
\end{algorithmic}
\end{algorithm}

\section*{Dry Run Example}
Consider the one‑dimensional quadratic 
\[
f(w)=(w-3)^{2},\qquad \nabla f(w)=2(w-3).
\]
Set $w_{0}=0$, $\alpha=0.1$, and run three iterations.

\begin{center}
\begin{tabular}{c|c|c|c|c}
$k$ & $w_{k}$ & $g_{k}= \nabla f(w_{k})$ & $G_{k}$ & $w_{k+1}=w_{k}-\frac{\alpha}{\sqrt{G_{k}}}g_{k}$\\ \hline
0 & $0$ & $2(0-3)=-6$ & $G_{1}=0+(-6)^{2}=36$ & $w_{1}=0-\frac{0.1}{\sqrt{36}}(-6)=0+0.1\cdot\frac{6}{6}=0.1$\\[4pt]
1 & $0.1$ & $2(0.1-3)=-5.8$ & $G_{2}=36+(-5.8)^{2}=36+33.64=69.64$ & $w_{2}=0.1-\frac{0.1}{\sqrt{69.64}}(-5.8)
      =0.1+0.1\frac{5.8}{8.34}\approx0.1695$\\[4pt]
2 & $0.1695$ & $2(0.1695-3)=-5.661$ & $G_{3}=69.64+(-5.661)^{2}=69.64+32.06=101.70$ &
 $w_{3}=0.1695-\frac{0.1}{\sqrt{101.70}}(-5.661)
   =0.1695+0.1\frac{5.661}{10.08}\approx0.2301$
\end{tabular}
\end{center}
Objective values: $f(w_{0})=9$, $f(w_{1})\approx8.41$, $f(w_{2})\approx8.02$, $f(w_{3})\approx7.66$.

\newpage
\section*{Assumptions}
\begin{assumption}[Convexity]\label{as:convex}
$f$ is convex, i.e.
\[
\forall u,v\in\mathbb{R}^{d}:\quad
f(v)\ge f(u)+\langle \nabla f(u),\,v-u\rangle .
\]
\end{assumption}
\begin{assumption}[Smoothness]\label{as:smooth}
$f$ is $L$‑smooth:
\[
\forall u,v\in\mathbb{R}^{d}:\quad
\|\nabla f(u)-\nabla f(v)\|\le L\|u-v\|.
\]
\end{assumption}
\begin{assumption}[Bounded Gradients]\label{as:bound}
There exists $G_{\max}>0$ such that $\|g_{k}\|=\|\nabla f(w_{k})\|\le G_{\max}$ for all $k$.
\end{assumption}
\begin{assumption}[Compact Feasible Set]\label{as:diam}
Let $w^{\star}\in\arg\min_{w}f(w)$. The diameter of the set $\{w_{0},w_{1},\dots\}$ is bounded:
\[
D:=\max_{k}\|w_{k}-w^{\star}\|<\infty .
\]
\end{assumption}

\section*{Main Convergence Theorem}
\begin{theorem}[AdaGrad Convergence for Convex Smooth Objectives]\label{thm:adagrad}
Under Assumptions \ref{as:convex}–\ref{as:diam}, after $T\ge 1$ iterations the average iterate
\[
\bar w_{T}:=\frac{1}{T}\sum_{k=1}^{T}w_{k}
\]
satisfies
\[
\boxed{
\mathbb{E}\bigl[f(\bar w_{T})-f(w^{\star})\bigr]
\;\le\;
\frac{D^{2}}{2\alpha}\,\frac{1}{\sum_{k=1}^{T}\frac{1}{\sqrt{G_{k}}}}
\;+\;
\frac{\alpha}{2T}\,\frac{\sum_{k=1}^{T}\|g_{k}\|^{2}}{\sum_{k=1}^{T}\frac{1}{\sqrt{G_{k}}}} }.
\]
In particular, using the bound $\|g_{k}\|\le G_{\max}$ and the monotonicity
$G_{k}\ge k\,G_{\max}^{2}/d$ (scalar case $d=1$) yields the simplified rate
\[
\mathbb{E}\bigl[f(\bar w_{T})-f(w^{\star})\bigr]
\;\le\;
\frac{DG_{\max}}{\sqrt{T}} .
\]
\end{theorem}

\section*{Theoretical Properties}
\begin{itemize}[leftmargin=*]
\item \textbf{Stability:} The denominator $\sqrt{G_{k}}$ grows monotonically, guaranteeing that the stepsizes $\eta_{k}$ are non‑increasing and thus preventing explosion.
\item \textbf{Hyperparameter Sensitivity:} The only free scalar is $\alpha$. The bound in Theorem~\ref{thm:adagrad} is minimized by $\alpha^{\star}=D/\sqrt{\sum_{k=1}^{T}\|g_{k}\|^{2}}$, reproducing the classic $O(DG_{\max}/\sqrt{T})$ rate.
\item \textbf{Comparison to SGD:} With a fixed stepsize $\eta$, SGD obtains $O(1/\sqrt{T})$ under the same assumptions. AdaGrad automatically adapts to the observed gradient magnitudes and enjoys the same asymptotic rate without manual tuning.
\item \textbf{Extension to Coordinates:} The scalar analysis extends component‑wise, yielding the standard AdaGrad bound $\sum_{j=1}^{d} D_{j}\sqrt{\sum_{k=1}^{T}g_{k,j}^{2}}$.
\end{itemize}

\section*{Discussion}
The theorem shows that the adaptive learning rate derived from the accumulated squared gradients yields a convergence guarantee identical to the optimal $O(1/\sqrt{T})$ rate for convex smooth problems, while automatically scaling with the observed gradient norm. The proof follows the classic online‑convex‑optimization (regret) argument, specialized to the deterministic gradient setting.

\newpage
\section*{Preliminaries (Definitions and Lemmas)}
\begin{lemma}[Convexity Inequality]\label{lem:conv}
For any $w$ and optimal $w^{\star}$,
\[
f(w)-f(w^{\star})\le\langle g, w-w^{\star}\rangle,\qquad g:=\nabla f(w).
\]
\end{lemma}
\begin{lemma}[Three‑Point Identity]\label{lem:3pt}
For any vectors $a,b,c\in\mathbb{R}^{d}$ and any scalar $\eta>0$,
\[
\|a-c\|^{2}= \|b-c\|^{2}+2\langle a-b,\,b-c\rangle+\|a-b\|^{2}.
\]
\end{lemma}
\begin{lemma}[Monotonicity of $G_{k}$]\label{lem:mon}
The sequence $G_{k}=\sum_{i=1}^{k}\|g_{i}\|^{2}$ is non‑decreasing and $G_{k}\ge\|g_{k}\|^{2}$.
\end{lemma}
\begin{proof}
Immediate from definition. \qedhere
\end{proof}

\section*{Main Proof (Step‑by‑Step Derivation)}
We prove Theorem~\ref{thm:adagrad}. All expectations are taken w.r.t.\ any randomness in the algorithm; in the deterministic setting they disappear.

\noindent\textbf{Step 1.} Expand the squared distance after an update using Lemma~\ref{lem:3pt} with $a=w_{k+1}$, $b=w_{k}$ and $c=w^{\star}$:
\[
\|w_{k+1}-w^{\star}\|^{2}
 =\|w_{k}-w^{\star}\|^{2}
   -2\eta_{k}\langle g_{k},\,w_{k}-w^{\star}\rangle
   +\eta_{k}^{2}\|g_{k}\|^{2}.
\tag{1}
\]

\noindent\textbf{Step 2.} Rearrange (1) to isolate the inner product term:
\[
2\eta_{k}\langle g_{k},\,w_{k}-w^{\star}\rangle
 =\|w_{k}-w^{\star}\|^{2}
   -\|w_{k+1}-w^{\star}\|^{2}
   +\eta_{k}^{2}\|g_{k}\|^{2}.
\tag{2}
\]

\noindent\textbf{Step 3.} Apply Lemma~\ref{lem:conv} to bound the suboptimality:
\[
f(w_{k})-f(w^{\star})\le\langle g_{k},\,w_{k}-w^{\star}\rangle.
\tag{3}
\]

\noindent\textbf{Step 4.} Multiply (3) by $2\eta_{k}$ and substitute the right‑hand side using (2):
\[
2\eta_{k}\bigl[f(w_{k})-f(w^{\star})\bigr]
 \le
 \|w_{k}-w^{\star}\|^{2}
 -\|w_{k+1}-w^{\star}\|^{2}
 +\eta_{k}^{2}\|g_{k}\|^{2}.
\tag{4}
\]

\noindent\textbf{Step 5.} Divide (4) by $2\eta_{k}$ and rearrange:
\[
f(w_{k})-f(w^{\star})
 \le
 \frac{\|w_{k}-w^{\star}\|^{2}-\|w_{k+1}-w^{\star}\|^{2}}{2\eta_{k}}
 +\frac{\eta_{k}}{2}\|g_{k}\|^{2}.
\tag{5}
\]

\noindent\textbf{Step 6.} Insert the explicit stepsize $\eta_{k}= \alpha/\sqrt{G_{k}}$ (note that $G_{k}$ is defined after the $k$‑th gradient is observed, i.e. $G_{k}= \sum_{i=1}^{k}\|g_{i}\|^{2}$):
\[
f(w_{k})-f(w^{\star})
 \le
 \frac{\sqrt{G_{k}}}{2\alpha}
   \bigl[\|w_{k}-w^{\star}\|^{2}-\|w_{k+1}-w^{\star}\|^{2}\bigr]
 +\frac{\alpha}{2\sqrt{G_{k}}}\|g_{k}\|^{2}.
\tag{6}
\]

\noindent\textbf{Step 7.} Sum inequality (6) over $k=1,\dots,T$:
\[
\sum_{k=1}^{T}\bigl[f(w_{k})-f(w^{\star})\bigr]
 \le
 \frac{1}{2\alpha}
   \sum_{k=1}^{T}\sqrt{G_{k}}
   \bigl[\|w_{k}-w^{\star}\|^{2}-\|w_{k+1}-w^{\star}\|^{2}\bigr]
 +\frac{\alpha}{2}
   \sum_{k=1}^{T}\frac{\|g_{k}\|^{2}}{\sqrt{G_{k}}}.
\tag{7}
\]

\noindent\textbf{Step 8.} The first sum on the right telescopes after weighting by $\sqrt{G_{k}}$. Using Lemma~\ref{lem:mon} we have $\sqrt{G_{k}}\le\sqrt{G_{T}}$ for all $k$, thus
\[
\sum_{k=1}^{T}\sqrt{G_{k}}
   \bigl[\|w_{k}-w^{\star}\|^{2}-\|w_{k+1}-w^{\star}\|^{2}\bigr]
 \le
 \sqrt{G_{T}}\bigl[\|w_{1}-w^{\star}\|^{2}-\|w_{T+1}-w^{\star}\|^{2}\bigr]
 \le
 \sqrt{G_{T}}\,D^{2},
\tag{8}
\]
where $D$ is the diameter defined in Assumption~\ref{as:diam}.

\noindent\textbf{Step 9.} Insert (8) into (7):
\[
\sum_{k=1}^{T}\bigl[f(w_{k})-f(w^{\star})\bigr]
 \le
 \frac{D^{2}\sqrt{G_{T}}}{2\alpha}
 +\frac{\alpha}{2}\sum_{k=1}^{T}\frac{\|g_{k}\|^{2}}{\sqrt{G_{k}}}.
\tag{9}
\]

\noindent\textbf{Step 10.} Divide both sides of (9) by $T$ and use the convexity of $f$ to replace the average of the function values by the function value at the average iterate $\bar w_{T}$:
\[
f(\bar w_{T})-f(w^{\star})
 \le
 \frac{1}{T}\sum_{k=1}^{T}\bigl[f(w_{k})-f(w^{\star})\bigr]
 \le
 \frac{D^{2}\sqrt{G_{T}}}{2\alpha T}
 +\frac{\alpha}{2T}\sum_{k=1}^{T}\frac{\|g_{k}\|^{2}}{\sqrt{G_{k}}}.
\tag{10}
\]

\noindent\textbf{Step 11.} Observe that $\sqrt{G_{k}}$ is non‑decreasing, therefore
\[
\frac{1}{\sqrt{G_{k}}}\le\frac{1}{\sqrt{G_{T}}},\qquad
\forall k\le T.
\]
Hence
\[
\sum_{k=1}^{T}\frac{\|g_{k}\|^{2}}{\sqrt{G_{k}}}
 \le
 \frac{1}{\sqrt{G_{T}}}\sum_{k=1}^{T}\|g_{k}\|^{2}
 =\frac{G_{T}}{\sqrt{G_{T}}}= \sqrt{G_{T}}.
\tag{11}
\]

\noindent\textbf{Step 12.} Substitute (11) into (10):
\[
f(\bar w_{T})-f(w^{\star})
 \le
 \frac{D^{2}\sqrt{G_{T}}}{2\alpha T}
 +\frac{\alpha}{2T}\sqrt{G_{T}}
 =\frac{\sqrt{G_{T}}}{2T}\Bigl(\frac{D^{2}}{\alpha}+\alpha\Bigr).
\tag{12}
\]

\noindent\textbf{Step 13.} Optimize the bound w.r.t.\ $\alpha$. The expression $\frac{D^{2}}{\alpha}+\alpha$ is minimized at $\alpha^{\star}=D$, giving the minimal value $2D$. Plugging $\alpha=D$ into (12) yields
\[
f(\bar w_{T})-f(w^{\star})
 \le
 \frac{D\sqrt{G_{T}}}{T}.
\tag{13}
\]

\noindent\textbf{Step 14.} Using the bounded‑gradient assumption $\|g_{k}\|\le G_{\max}$ we have $G_{T}\le T G_{\max}^{2}$, thus $\sqrt{G_{T}}\le G_{\max}\sqrt{T}$. Inserting into (13) gives the final $O(1/\sqrt{T})$ rate:
\[
f(\bar w_{T})-f(w^{\star})
 \le
 \frac{DG_{\max}}{\sqrt{T}}.
\tag{14}
\]

All steps above are reversible under the stated assumptions; therefore inequality (14) holds for every $T\ge1$, completing the proof. \qed

\section*{Rate Derivation (With Constants)}
From (12) we have the explicit bound
\[
f(\bar w_{T})-f(w^{\star})
 \le
 \underbrace{\frac{D^{2}}{2\alpha}}_{C_{1}}
 \frac{\sqrt{G_{T}}}{T}
 +\underbrace{\frac{\alpha}{2}}_{C_{2}}
 \frac{\sqrt{G_{T}}}{T}.
\]
Choosing $\alpha=D$ balances the two constants $C_{1}=C_{2}=D/2$, and the bound becomes
\[
f(\bar w_{T})-f(w^{\star})\le\frac{D\sqrt{G_{T}}}{T}.
\]
If additionally $G_{T}\le T G_{\max}^{2}$,
\[
f(\bar w_{T})-f(w^{\star})\le\frac{DG_{\max}}{\sqrt{T}}.
\]

\section*{Special Cases}
\begin{enumerate}[leftmargin=*]
\item \textbf{Exact Quadratic $f(w)=\tfrac{1}{2}w^{\top}Qw+b^{\top}w$ with $Q\succeq 0$.}  
Because gradients are linear, $g_{k}=Qw_{k}+b$ and $G_{k}$ grows at most linearly in $k$, yielding the same $O(1/\sqrt{T})$ bound; moreover, if $Q=0$ (pure linear function) the algorithm converges in a single step because $G_{k}$ becomes constant after the first iteration.

\item \textbf{Strongly Convex $f$.}  
If $f$ is $\mu$‑strongly convex, a refined analysis (adding the strong‑convexity inequality) gives a faster $O\bigl((\log T)/T\bigr)$ rate; the proof proceeds by replacing Lemma~\ref{lem:conv} with the strong‑convexity bound.

\item \textbf{Coordinate‑wise AdaGrad.}  
Applying the scalar proof component‑wise yields the classic bound
\[
\sum_{j=1}^{d} D_{j}\sqrt{\sum_{k=1}^{T} g_{k,j}^{2}},
\]
which reduces to the scalar case when $d=1$.
\end{enumerate}

\end{document}